\chapter{Retribución en CaixaBank}

\section{Análisis de puestos representativos y consulta del Convenio Colectivo sectorial}

La política retributiva de CaixaBank se enmarca en el \textbf{IV Convenio Colectivo del Sector de Banca}, negociado entre la Asociación Española de Banca (AEB) y los principales sindicatos del sector (CCOO, UGT y FINE). Este convenio establece las categorías profesionales, las tablas salariales mínimas y las condiciones laborales de aplicación general para todas las entidades bancarias adheridas, actuando como suelo retributivo que CaixaBank respeta y, en la mayoría de los casos, supera.

Con el fin de reflejar la diversidad de funciones y niveles jerárquicos de la entidad, se han seleccionado cinco puestos representativos distribuidos a lo largo de su estructura organizativa, desde los perfiles operativos hasta la alta dirección:

\begin{itemize}
    \item \textbf{Operativo de caja / Auxiliar administrativo (Nivel I):} Es el puesto de entrada en la estructura de sucursal. Sus funciones abarcan la atención básica al cliente en ventanilla, la gestión de efectivo y las operaciones de rutina. Según el convenio colectivo sectorial vigente, el salario base mínimo para esta categoría se sitúa en torno a los \textbf{17.500\,€ brutos anuales}, aunque en CaixaBank la retribución total efectiva, incluyendo complementos de antigüedad y pluses de convenio, alcanza habitualmente los \textbf{20.000--22.000\,€ brutos anuales}.

    \item \textbf{Gestor de Banca Personal (Nivel II/III):} Perfil comercial encargado de asesorar a clientes particulares en productos de ahorro, inversión básica y financiación. Requiere la certificación MiFID II. El convenio fija su retribución mínima en aproximadamente \textbf{22.000\,€ brutos anuales}; en CaixaBank, el paquete retributivo total —incluyendo variable comercial— oscila entre los \textbf{28.000\,€ y los 35.000\,€ brutos anuales} en función de los objetivos alcanzados.

    \item \textbf{Gestor de Banca Privada / Analista Senior (Nivel III/IV):} Especialista en asesoramiento patrimonial a clientes de alto patrimonio neto. Exige la certificación EFPA o EFA y habilidades analíticas avanzadas. Es uno de los perfiles más valorados del sector, con una retribución que en CaixaBank se aproxima a los \textbf{45.000--60.000\,€ brutos anuales}, a lo que se añade una parte variable significativa vinculada al volumen de patrimonio gestionado.

    \item \textbf{Director/a de Sucursal (Nivel IV/V):} Responsable de la gestión integral de una oficina bancaria: liderazgo del equipo, captación de negocio, cumplimiento normativo y relación con clientes clave. Su retribución en CaixaBank ronda los \textbf{50.000--70.000\,€ brutos anuales}, incluyendo un variable ligado a los resultados globales de la sucursal.

    \item \textbf{Director/a de Área o Ejecutivo/a Senior (Nivel V/Directivo):} Perfiles de alta responsabilidad que supervisan varias sucursales o lideran departamentos corporativos. Su retribución total puede superar los \textbf{100.000\,€ brutos anuales}, complementada con incentivos a largo plazo como acciones del banco o planes de co-inversión.
\end{itemize}

Cabe destacar que CaixaBank cotiza al Ibex 35 y publica anualmente su brecha salarial de género y sus ratios retributivos en el informe de sostenibilidad, donde se recoge que la retribución media de su plantilla supera en un 30--40\,\% los mínimos establecidos por el convenio sectorial, lo que refleja una apuesta decidida por la competitividad retributiva como palanca de atracción y retención del talento.

\section{Gestión de la retribución: los tres tipos de equidad}

Una gestión retributiva eficaz debe garantizar que los empleados perciban su salario como justo, tanto en relación con sus compañeros como con el mercado externo y con su propia aportación individual. CaixaBank articula su política de compensación en torno a los tres tipos clásicos de equidad retributiva.

\subsection{Equidad interna}

La equidad interna hace referencia a la coherencia salarial entre los distintos puestos de trabajo dentro de la propia organización. Para garantizarla, CaixaBank estructura su plantilla en un sistema de \textbf{niveles o bandas salariales} —denominados internamente grupos profesionales— que clasifica todos los puestos según su valor relativo para la organización, atendiendo a criterios como la responsabilidad asumida, la complejidad de las tareas, el impacto en el negocio y las competencias técnicas requeridas.

Este modelo asegura que dos empleados que desempeñen funciones de igual valor y complejidad no experimenten diferencias salariales injustificadas. La implantación de \textbf{SAP SuccessFactors} como plataforma central de gestión de personas facilita la administración de estas bandas y permite a los responsables de Recursos Humanos detectar desviaciones internas con rapidez, lo que contribuye a mantener la coherencia de la estructura retributiva en los más de 45.000 empleados del grupo.

\subsection{Equidad externa}

La equidad externa mide la competitividad de la retribución de la empresa respecto a otras organizaciones del sector o del mercado de referencia. Para garantizarla, el Área de Personas y Organización de CaixaBank realiza anualmente \textbf{estudios de benchmarking salarial} apoyándose en las encuestas de compensación de consultoras especializadas como Mercer o Willis Towers Watson, que recogen los datos del sector financiero español e internacional.

Estos estudios permiten a la entidad conocer en qué posición se sitúa su paquete retributivo en comparación con sus competidores directos —BBVA, Banco Santander, Sabadell o Bankinter—, y ajustar sus tablas salariales para permanecer en una posición competitiva. CaixaBank tiene como objetivo declarado situarse en el cuartil superior del mercado bancario español para los perfiles más críticos, con el fin de atraer y retener el talento especializado que el proceso de transformación digital del grupo exige.

\subsection{Equidad individual}

La equidad individual garantiza que, dentro de una misma banda salarial, se reconozcan las diferencias de rendimiento y aportación de cada empleado. CaixaBank materializa este principio a través de un sistema de \textbf{retribución variable ligada al desempeño}, cuya base es la evaluación de rendimiento anual gestionada mediante SAP SuccessFactors.

En función de los resultados obtenidos, el rendimiento de cada empleado se clasifica en distintos niveles de desempeño que determinan el porcentaje de retribución variable al que tiene derecho. De esta forma, un gestor comercial que supere de manera significativa sus objetivos anuales percibirá un bonus considerablemente superior al de un compañero del mismo nivel que haya alcanzado únicamente los mínimos. Este mecanismo busca no solo reconocer el mérito individual, sino también mantener el incentivo y la motivación de los equipos de alto rendimiento.

\section{Retribución por puesto vs.\ por competencias}

Uno de los debates centrales en el diseño de sistemas retributivos es si la base de la compensación debe ser el puesto de trabajo en sí mismo —con independencia de quién lo ocupe— o bien las competencias y capacidades que el empleado aporta a la organización. CaixaBank aplica un \textbf{modelo híbrido} que combina ambos enfoques de forma complementaria.

Por un lado, la \textbf{retribución basada en el puesto} constituye el componente fijo y estable del salario. El convenio colectivo sectorial y la estructura de niveles profesionales internos de la entidad asignan a cada categoría un rango salarial de referencia, determinado por las responsabilidades formales del puesto, su posición en el organigrama y el impacto esperado en los resultados del negocio. Esta base fija garantiza la equidad interna y la seguridad económica del empleado con independencia de las fluctuaciones del rendimiento en un periodo concreto.

Por otro lado, la \textbf{retribución basada en competencias} actúa como elemento diferenciador sobre esa base fija. En CaixaBank, determinadas competencias tienen un reconocimiento retributivo explícito:

\begin{itemize}
    \item La obtención y mantenimiento de \textbf{certificaciones financieras reguladas} —como la certificación MiFID II para gestores de banca personal o la titulación EFPA/EFA para gestores de banca privada— está directamente vinculada a la posibilidad de acceder a puestos de mayor categoría y, por ende, a bandas salariales superiores.
    \item Las \textbf{habilidades digitales avanzadas} —manejo de herramientas de análisis de datos, uso de inteligencia artificial generativa o gestión de proyectos tecnológicos— son cada vez más valoradas en el proceso de promoción interna, especialmente en el contexto de la transformación digital recogida en el Plan Estratégico 2025--2027.
    \item La \textbf{capacidad de liderazgo} y los resultados obtenidos en programas de desarrollo directivo internos influyen en la posibilidad de acceder a complementos salariales reservados a los perfiles de alta potencial (\textit{high potential}).
\end{itemize}

Este modelo mixto responde a la necesidad de combinar la estabilidad y la equidad que ofrece la retribución por puesto con la flexibilidad y el reconocimiento del mérito individual que aporta el enfoque por competencias, lo que resulta especialmente adecuado para una organización tan heterogénea como CaixaBank.

\section{Tipos de incentivos y elementos no financieros}

La política de compensación total de CaixaBank va mucho más allá del salario base. La entidad diseña un paquete retributivo integral que combina incentivos financieros directos e indirectos con un conjunto de beneficios no financieros cuyo objetivo es mejorar la calidad de vida de sus empleados y fortalecer su compromiso con la organización.

\subsection{Incentivos financieros directos: retribución variable}

\begin{itemize}
    \item \textbf{Bonus comercial anual:} Es el principal incentivo financiero para los perfiles de red de oficinas. Se calcula como un porcentaje del salario fijo y se determina en función del grado de cumplimiento de los objetivos individuales y colectivos (captación de clientes, volumen de contratación de productos, calidad del servicio). En años de buen rendimiento, este componente puede suponer entre el 10\,\% y el 30\,\% de la retribución total del gestor.
    \item \textbf{Planes de co-inversión y acciones para directivos:} Para los perfiles de alta dirección, CaixaBank dispone de planes de incentivos a largo plazo basados en la entrega diferida de acciones del banco. Estos programas vinculan parte de la retribución al desempeño del grupo a lo largo de tres o cinco años, alineando los intereses de los directivos con los de los accionistas y favoreciendo la retención del talento directivo.
    \item \textbf{Incentivos por captación y fidelización:} Algunos programas retributivos contemplan primas puntuales por la incorporación de clientes de alto valor o por la renovación de contratos de productos de largo plazo.
\end{itemize}

\subsection{Retribución indirecta y prestaciones}

\begin{itemize}
    \item \textbf{Seguro médico colectivo:} CaixaBank ofrece a todos sus empleados la posibilidad de acceder a un seguro de salud de empresa con condiciones ventajosas, extensible a los familiares directos. Esta prestación tiene un valor percibido muy alto y contribuye significativamente a la satisfacción del empleado.
    \item \textbf{Plan de pensiones de empresa:} La entidad realiza aportaciones periódicas al plan de pensiones de sus empleados como complemento a la jubilación pública, lo que constituye uno de los elementos retributivos indirectos más apreciados, especialmente entre los empleados de mayor antigüedad.
    \item \textbf{Préstamos hipotecarios y personales a tipo preferencial:} Los trabajadores de CaixaBank tienen acceso a productos de financiación con condiciones más favorables que las del mercado, lo que supone un beneficio económico tangible y directo.
    \item \textbf{Ayuda escolar y guardería:} La entidad contempla subsidios o ayudas económicas para el cuidado de hijos menores, en línea con su compromiso con la conciliación de la vida laboral y familiar.
\end{itemize}

\subsection{Elementos no financieros}

Los elementos no financieros de la retribución constituyen un componente esencial del paquete de compensación total de CaixaBank, especialmente relevantes para atraer a las generaciones más jóvenes:

\begin{itemize}
    \item \textbf{Formación y desarrollo profesional:} El acceso a la plataforma Virtaula, a los programas de las escuelas especializadas del banco y la financiación de certificaciones externas (EFPA, CFA, etc.) representa una inversión en el capital humano del empleado que este percibe como parte de su retribución total.
    \item \textbf{Flexibilidad horaria:} CaixaBank dispone de medidas de flexibilidad en la entrada y salida, así como de jornada intensiva en el periodo estival, lo que mejora la conciliación laboral.
    \item \textbf{Reconocimiento y clima laboral:} La entidad dispone de programas de reconocimiento interno —como premios a la innovación o a la excelencia comercial— que refuerzan el sentido de pertenencia y la motivación sin implicar necesariamente un coste económico directo.
    \item \textbf{Voluntariado corporativo:} La vinculación con la Fundación ``la Caixa'' permite a los empleados participar en actividades de responsabilidad social, lo que refuerza el orgullo de pertenencia y la identificación con los valores de la entidad.
\end{itemize}

\section{Tendencias: retribución flexible y salario emocional}

La gestión de la retribución está evolucionando hacia modelos más personalizados y centrados en el bienestar integral del empleado. CaixaBank ha incorporado dos de las tendencias más relevantes en este ámbito: la retribución flexible y el salario emocional.

\subsection{Retribución flexible}

La retribución flexible —también denominada \textit{flex benefits} o salario en especie— consiste en permitir que el empleado decida voluntariamente transformar una parte de su retribución dineraria en una serie de productos o servicios exentos total o parcialmente de tributación en el Impuesto sobre la Renta de las Personas Físicas (IRPF), optimizando así su poder adquisitivo neto sin incrementar el coste laboral para la empresa.

CaixaBank ofrece a sus empleados la posibilidad de acogerse a un plan de retribución flexible que incluye, entre otros, los siguientes productos:

\begin{itemize}
    \item \textbf{Seguro médico:} Los empleados pueden destinar parte de su salario bruto al pago de la prima del seguro médico privado, beneficiándose de la exención fiscal establecida en el artículo 42 de la Ley del IRPF.
    \item \textbf{Transporte:} Bonificaciones para el uso de transporte público colectivo, con el correspondiente ahorro fiscal.
    \item \textbf{Ticket restaurante:} Vales de comida para los días de trabajo presencial, exentos de IRPF hasta los límites legales.
    \item \textbf{Guardería y cheque escolar:} Contribución a los gastos de cuidado de hijos menores de tres años, con beneficio fiscal para el trabajador.
    \item \textbf{Formación:} Determinados programas de formación financiados por la empresa están excluidos de la base imponible del impuesto.
\end{itemize}

Este sistema aporta a los empleados un aumento efectivo de su renta disponible sin modificar el coste total de personal para la entidad, lo que lo convierte en una herramienta muy eficiente para incrementar la satisfacción retributiva. Además, permite a cada trabajador personalizar su paquete de beneficios en función de su situación personal y familiar, respondiendo a la creciente demanda de individualización de las condiciones laborales.

\subsection{Salario emocional}

El concepto de salario emocional hace referencia al conjunto de beneficios no económicos que una organización ofrece a sus empleados y que tienen un impacto directo en su satisfacción, motivación y bienestar, sin que su valor se traduzca directamente en un ingreso monetario. En un contexto de mercado de talento competitivo, el salario emocional se ha convertido en un factor diferenciador decisivo para la atracción y retención de profesionales, especialmente entre las generaciones más jóvenes.

CaixaBank ha desarrollado diversas iniciativas en esta línea, agrupadas bajo su programa de bienestar y conciliación:

\begin{itemize}
    \item \textbf{Teletrabajo estructurado:} Tras la aceleración impuesta por la pandemia, la entidad ha institucionalizado el trabajo en remoto para los puestos elegibles, permitiendo hasta dos días de teletrabajo semanales. Esta medida impacta directamente en la reducción de los tiempos de desplazamiento y en la autonomía del empleado para organizar su jornada.
    \item \textbf{Medidas de conciliación:} CaixaBank dispone de una batería de medidas que incluyen días adicionales de asuntos propios, permisos de maternidad y paternidad ampliados voluntariamente más allá de los mínimos legales, y adaptaciones de jornada para el cuidado de familiares dependientes.
    \item \textbf{Bienestar mental y físico:} La entidad ha ampliado su oferta de bienestar con programas de apoyo psicológico, acceso a plataformas de mindfulness y actividad física, y talleres de gestión del estrés, en respuesta al incremento de las bajas por salud mental en el sector bancario.
    \item \textbf{Voluntariado y propósito:} La participación en los programas sociales de la Fundación ``la Caixa'' —que opera en ámbitos como la inserción laboral, la investigación médica o la educación— proporciona a los empleados un sentido de propósito que trasciende su función laboral diaria, aumentando su vinculación emocional con la organización.
    \item \textbf{Desarrollo y formación como beneficio percibido:} El acceso privilegiado a un ecosistema de formación continua —desde las escuelas internas hasta las certificaciones financieras costeadas por el banco— es percibido por los empleados como una inversión en su futuro profesional que difícilmente encontrarían en otra entidad, lo que actúa como un poderoso elemento de retención.
\end{itemize}

En definitiva, la combinación de retribución flexible y salario emocional responde a la comprensión de que la motivación del empleado moderno no depende únicamente de su nivel salarial, sino de la calidad de su experiencia global como trabajador. CaixaBank, consciente de ello, integra ambas tendencias en su propuesta de valor al empleado (\textit{Employee Value Proposition}) como palanca estratégica para mantener una plantilla comprometida, productiva y alineada con los objetivos del grupo.
